\documentclass[11pt,a4paper]{article}
\usepackage[T1]{fontenc}
\usepackage{lmodern,microtype}
\usepackage[a4paper,margin=25mm]{geometry}
\usepackage{amsmath,amssymb,booktabs,array,tabularx}
\usepackage[dvipsnames]{xcolor}
\usepackage[colorlinks=true,linkcolor=MidnightBlue,urlcolor=MidnightBlue]{hyperref}
\usepackage{fancyhdr}
\pagestyle{fancy}
\fancyhf{}
\fancyhead[L]{\small Rank-five fusion categories}
\fancyhead[R]{\small Third revision round}
\fancyfoot[C]{\thepage}
\setlength{\headheight}{14pt}
\setlength{\emergencystretch}{2em}
\setlength{\parskip}{3pt}
\newcommand{\FPdim}{\operatorname{FPdim}}
\newcommand{\Tr}{\operatorname{Tr}}
\hypersetup{pdftitle={Rank-five classification: third revision round},pdfauthor={}}
\title{Rank-five classification\\[4pt]\large Further audit and revision letter\\Third revision round}
\author{}
\date{24 November 2026}
\begin{document}
\maketitle

\section{Outcome and exact baseline}

This letter records all changes relative to the previously delivered
\path{Rank5_General_Sharp_Revised_Round2.zip}. The manuscript, its
computational proofs, and the previous repairs were reviewed again.
Two substantive defects were found and repaired:
\begin{enumerate}
\item The last conductor-five exclusion assumed that certain third
Frobenius--Schur indicators were real. This omitted possible
half-integer real parts. The corrected argument covers 600 indicator
tuples instead of 180. The existing six separating forms exclude all
600, closing the gap.
\item A subring-certificate checker accepted negative basis indices.
A concrete malformed certificate could therefore falsely exclude the
realized ring $\operatorname{Rep}(D_{14})$, where $D_{14}$ has order
fourteen. Strict index validation and a regression check repair this
defect. Neither of the two stored subring certificates was malformed.
\end{enumerate}
There is also one punctuation correction and a new independent check
of the exact dimension cutoff. The complete source and package changes
are listed below and in \path{revision_round3.patch}.

The corrected argument retains every principal conclusion: every
complex rank-five fusion category has a pseudo-unitary Galois conjugate
and admits a spherical structure, with
\[
\mu\le4,\qquad \FPdim(\mathcal C)\le48+24\sqrt3.
\]
There are exactly seventeen based rings, each with a unitary realization.
The 269-ring generation, 114-ring bounded census, 97 exclusions, and
17 model identifications are unchanged. The earlier reports remain
historical records; their acceptance of the former third-indicator
range is superseded by the correction in this letter.

\section{R3-M01: third indicators and the final conductor-five exclusion}

\textbf{Location:} Proposition 11.10, Section 11.5, the last
conductor-five case; source label \path{prop:sharp-exchange}.
The corresponding program is \path{checks/verify_exchange.py}.

\paragraph{The mistake.}
The previous proof asserted that duality pairs complex conjugate
higher indicators, and used this to regard
$\nu_3(X)+\nu_3(Y)$ and $\nu_3(Z)$ as real traces. In a spherical
category the relevant duality identity is instead
\[
                         \nu_n(V^*)=\nu_n(V).
\]
This follows by specializing Ng--Schauenburg,
\href{https://arxiv.org/pdf/0806.2493v3}{arXiv:0806.2493v3},
Proposition 3.3, to the unit central object. It does not force reality.
For a concrete illustration that higher indicators can be nonreal,
Ng--Schopieray--Wang,
\href{https://arxiv.org/pdf/1812.11234v2}{arXiv:1812.11234v2},
Example 4.6, records the regular third indicator
$3(5+4\zeta_3^2)$ for a semisimple Hopf algebra of dimension 27.

\paragraph{The complete real-part range.}
If a rotation of order dividing three acts on an $H$-dimensional
space and its eigenvalue-one multiplicity is $r$, then
\[
 \operatorname{Re}\Tr(T)
 =r-\frac{H-r}{2}=\frac{3r-H}{2},\qquad r=0,\ldots,H.
\]
In the case at issue,
$[X^3:1]+[Y^3:1]=s=4$ and $[Z^3:1]=e=19$.
The first equality holds whether $X,Y$ are self-dual or form a dual
pair: $X^2=Y^2$, and duality permutes $X,Y$.
Consequently the proof now uses
\[
\begin{gathered}
 a=\operatorname{Re}(\nu_3(X)+\nu_3(Y)),\qquad
 b=\operatorname{Re}\nu_3(Z),\\
 a\in\left\{\frac{3r-4}{2}:r=0,\ldots,4\right\},\qquad
 b\in\left\{\frac{3r-19}{2}:r=0,\ldots,19\right\}.
\end{gathered}
\]
There are five choices for $a$ and twenty for $b$. Together with
the six second-indicator choices, this gives $6\cdot5\cdot20=600$
tuples. The previous three and ten choices covered only 180;
\textbf{420 previously untested combinations are now included}.

\paragraph{Why the averaged moment argument still applies.}
The revised proof explicitly writes
\[
 m_j=\operatorname{Re}M_j,\qquad
 M_j=\sum_W\left(\frac{\dim W}{D}\right)^2\theta_W^j.
\]
The signed modular Galois permutation gives
$\sigma_k(M_j)=M_{j/k}$, hence
$\sigma_k(m_j)=m_{j/k}$, for $k$ coprime to the twist exponent.
The moments used belong to $K=\mathbb Q(\sqrt5)$: the first two
already lie in $K$, and
\[
                         m_3=p(1+ax+bz)\in K
\]
because real parts of cubic rotation traces are rational.
Averaging $\operatorname{Re}(\theta^j)$ over the two quadratic
Galois cosets therefore integrates to $m_j$ and $m_j'$;
inversion preserves each coset. This does not require $M_3$ itself
to be real.

\paragraph{Exact computational repair.}
The program now constructs both full ranges with rational fractions
and checks that there are 6 targets at $f=20$ and 600 at $f=5$.
All ten and twenty-five root columns, respectively, are still checked.
The same two and six integral separating forms remain nonnegative
on every root column and have the required signs on the mass-bound
coordinates. For every target, at least one of the relevant forms
has a strictly negative value.
No separator coefficient or twist-mass bound changed.

The earlier third-indicator congruence in Section 11 is unaffected:
there the temporary assumption $3\nmid N$ forces the indicators
into $\mathbb Q(\zeta_3)\cap\mathbb Q(\zeta_N)=\mathbb Q$.
The third-indicator arguments in Sections 15 and 16 use absolute
values or real parts and do not make the erroneous reality inference.

\section{R3-C01: invalid subring indices could create a false exclusion}

\textbf{Location:} \path{final_census/criteria/decisions.py},
function \texttt{check\_subring}.

The old checker required distinct entries in a subset, but did not
first verify that every entry was a valid basis index. In a rank-five
Python list, index $-4$ refers to the same entry as index $1$.
For the realized ring $\operatorname{Rep}(D_{14})$, the old code
accepted the certificate
\begin{verbatim}
{"kind": "subring", "subset": [0, 1, -4],
 "quotient_polynomial": ["1", "-14/3"]}
\end{verbatim}
The unit and nonidentity invertible have indices zero and one. The
malformed list counted the nonidentity invertible twice, assigned
the fictitious subring dimension three, and accepted the
nonintegrality of $14/3$ as an exclusion. Distinctness of the raw
indices did not detect their aliasing.

Before any tensor or character indexing, the checker now requires
\[
 \operatorname{type}(i)=\texttt{int},\qquad
                         0\le i<\operatorname{rank}(R)
\]
for every listed index. The existing distinctness, unit, size,
duality-closure, multiplication-closure, and quotient-polynomial
checks remain in place. The repaired checker rejects the concrete
certificate above.

The optional program \path{final_census/verify_additional_audits.py}
now constructs this realized model and requires rejection of that
exact certificate. Both stored subring certificates contain valid
indices and still verify; the mathematical subcategory-divisibility
criterion and its two actual exclusions are unchanged.

\section{R3-V01: independent comparison at the sharp dimension cutoff}

The optional additional audit now also checks the dimension decision
on all 269 generated rings. For each Casimir characteristic polynomial,
it isolates the largest root by exact rational intervals and compares
it with $M=48+24\sqrt3$. The endpoint's minimal polynomial is
\[
                         t^2-96t+576.
\]
Polynomial divisibility detects equality when the largest-root
interval contains $M$; otherwise rational interval refinement
separates the root from $M$. This calculation does not use the
shifted-coefficient criterion to make its own decision. It then
compares the two answers.

All 269 decisions agree: 114 rings satisfy the bound, and three of
the generated rings lie exactly at the endpoint. The latter count
concerns generated fusion rings before the categorification tests;
it does not assert three categorified dimension maximizers.
The original generator, dimension-test implementation, and threshold
are unchanged. The added audit compiles temporary binaries, uses
source paths relative to its own file, and has explicit
exception-based coverage checks.

\section{Verification and mathematical review}

All three commands passed after the applicable edits:
\begin{verbatim}
python3 verify.py
python3 final_census/verify_independent_census.py
python3 final_census/verify_additional_audits.py
\end{verbatim}
The ordinary run passed all eighteen structural checks, including
the enlarged 600-tuple exclusion, and reproduced the census:
\begin{center}
\begin{tabular}{rrrrr}
\toprule
$\mu$&Generated&Within both bounds&Excluded&Realized\\\midrule
1&16&16&6&10\\
2&37&30&27&3\\
3&82&46&44&2\\
4&134&22&20&2\\\midrule
Total&269&114&97&17\\\bottomrule
\end{tabular}
\end{center}
The exclusions are still 76 $d$-number failures, two subcategory
divisibility failures, three central-dimension norm failures, two
dimension-field failures, ten induction/second-moment failures,
and four higher-indicator certificates. All seventeen model
identifications pass.

The independent census replay again gives 220 self-dual rings,
47 rings with one nontrivial dual pair, and two with two dual pairs.
The independent Pieri construction agrees with the
$\operatorname{Ad}(\mathrm{SU}(3)_4)$ model.
The additional audit also passes its previous checks:
\begin{itemize}
\item the quadratic minimum $5/18$ and contradiction margin $509/16200$;
\item 864 exact cyclotomic moments and 19 joint trace vectors;
\item 14 induction/indicator certificate rings and 11 residual Gram matrices;
\item all three quartic generator enclosures and every reported upper endpoint.
\end{itemize}

The review reread all seventeen sections. The mathematical checks
included the sphericalization and common-field reductions, simultaneous
columns, Hecke/Fourier bounds, all quadratic/quartic/quintic descent
ranges, the invertible-action cases, $S_3$ untwisting, Morita induction,
conductor-eight saturation, the small-generator graph, quintic
maximal-order argument, local socles and indicator limits, twist-support
applications, cubic/quartic depth arguments, and the endpoint
equivariantization. Primary sources were checked where needed.

The new dimension comparison shares the existing ring generator;
the previous alternative census shares its linear-completion routine
and basis parametrization. The trace, Gram, and enclosure checks
share stated field data, root intervals, or individual column
construction. These limitations are now reflected in both READMEs.
No further confirmed error was found. This audit does not claim to
reprove every cited theorem or to give a formal proof-assistant
certification.

\section{Complete remaining change ledger}

\begin{enumerate}
\item \textbf{R3-M02, punctuation.} In Proposition 11.1, removed
the stray comma between ``the nontrivial quadratic character'' and
``gives''. There is no mathematical change in that sentence.
\item \textbf{Manuscript source and PDF.} The source contains R3-M01
and R3-M02, and the PDF was rebuilt. The title, empty author fields,
theorem numbering, principal bounds, and seventeen-ring list are retained.
\item \textbf{Code.} The only edited code files are
\path{checks/verify_exchange.py},
\path{final_census/criteria/decisions.py}, and
\path{final_census/verify_additional_audits.py}, with precisely the
range, validation, regression, and cutoff-audit changes described above.
The last file's imports, module description, and main routine were
updated for its two added checks.
\item \textbf{Documentation.} The root README identifies the current
letter, patch, and transcript, preserves the status of earlier reports,
documents the 600-tuple repair, and describes the new optional checks.
The census README explains the corrected index validation and the
independent dimension comparison, and updates the file-role and
transcript references. Both state the shared dependencies.
\item \textbf{Recorded outputs.} \path{verification.txt} is the new
ordinary run; its changed endpoint line records 600 targets.
\path{additional_verification_round3.txt} is new and contains the
actual outputs of the two optional commands with command labels.
The preceding independent and second-round transcripts are preserved.
\item \textbf{New letter and exact diff.} This letter is new in
\path{Revision_Letter_Round3.tex} and
\path{Revision_Letter_Round3.pdf}.
\path{revision_round3.patch} is new and records every edit to an
existing text file, apart from the regenerated checksum manifest.
The three earlier report pairs and both earlier patches are preserved.
\item \textbf{Final package.} \path{SHA256SUMS.txt} was regenerated;
every distributed checksum and the archive were checked. Temporary
build files, compiled helpers, and Python caches are excluded. The
final manuscript and letter were rendered and visually inspected.
\end{enumerate}
No stored witness JSON, model tensor, generator source, separator
coefficient, or dependency requirement changed.

For exact traceability, the preceding delivered ZIP has SHA-256
\begin{center}\scriptsize\ttfamily
0981d34ee027e572c1933df4cdcf978884e207b1c5977cd6fac675e9743e3054
\end{center}
and its \path{rank5.tex} has SHA-256
\begin{center}\scriptsize\ttfamily
3837538e5e63736a62bb13d3a63d5f1f032c472dd68a1d57058cc00e561047c2
\end{center}
Every change in this letter is measured against that exact revision.

\end{document}
